\thispagestyle{empty}
\begin{center}
\begin{LARGE}
\textbf{Resumen}
\end{LARGE}
\end{center}
\begin{quotation}
\sloppy

Se presenta el desarrollo y la evaluación experimental de un prototipo de sistema automatizado de registro de asistencia mediante reconocimiento facial biométrico en la Facultad Politécnica de la Universidad Nacional del Este (FPUNE), Paraguay. El método manual actual consume entre 5 y 7 minutos por sesión y presenta tasas de error del 15–20\% debido a registros ilegibles y omisiones. Mediante un diseño experimental con \textit{N}=13 estudiantes de la institución durante 8 semanas, se implementó un sistema de dos etapas: la primera siendo la detección facial con \textit{YOLOv8n-Face} y la segunda correspondiente a la identificación mediante embeddings \textit{ArcFace}. Se recopilaron 65 imágenes por estudiante bajo condiciones variables de iluminación y distancias comprendidas entre 1 a 5m, evaluadas en 24 sesiones lectivas reales con un total de 1080 registros. Los resultados mostraton una precisión promedio del 84\% destacando la iluminación como factor determinante en el rendimiento, con tiempos de procesamiento de 288±23\,ms por rostro. La implementación redujo el tiempo de registro, de unos 5min promedio a unos 1.35min y eliminó los errores de transcripción manual, aunque el rendimiento depende críticamente de la iluminación. Los resultados validan la viabilidad técnica y económica del sistema propuesto y demuestran su potencial para mejorar la gestión académica y administrativa en instituciones educativas.

\vspace*{0.5cm}
\noindent \textbf{Descriptores:} 1. Reconocimiento facial, 2. Biometría, 3. Aprendizaje profundo, 4. Asistencia automatizada, 5. Visión por computadora, 6. Gestión académica.

\end{quotation}
